\subsection{可迁移关系}

让 \(\mathfrak{C}\) 是一个比集合论更强的理论，设 \(x_{1},\ldots ,x_{n},\) \(s_1,\ldots ,s_p\) 是互不相同的字母，且与 \(\mathfrak{C}\) 的常元不同，并且设 \(\mathrm{A}_1,\ldots ,\mathrm{A}_m\) 为 \(\mathfrak{C}\) 中的项，其中字母 \(x_{i}\) \((1\leqslant i\leqslant n)\) 与 \(s_j\) \((1\leqslant j\leqslant p)\) 均不出现。设 \(\mathrm{S}_1,\ldots ,\mathrm{S}_p\) 是 \(n + m\) 项上的阶梯构造方案。则关系 \(\mathrm{T}\{x_1,\dots,x_n,s_1,\dots,s_p\}\) :

\[
\begin{array}{c}
\text{``}s_{1} \in S_{1}(x_{1}, \ldots, x_{n}, A_{1}, \ldots, A_{m})
\text{ 且 }s_{2} \in S_{2}(x_{1}, \ldots, x_{n}, A_{1}, \ldots, A_{m}) \\
\text{且 }\ldots\text{ 且 }s_{p} \in S_{p}(x_{1}, \ldots, x_{n}, A_{1}, \ldots, A_{m})\text{''}
\end{array}
\]

被称为字母的典型化 \(s_{1}, \ldots, s_{p}\) .

设 \(\mathbb{R}\{x_1, \ldots, x_n, s_1, \ldots, s_p\}\) 是 \(\mathfrak{C}\) 中的一个关系，其中含有某些字母 \(x_i, s_j\) （也可能含有其他字母）。如果满足下述条件，则称 \(\mathbb{R}\) 在 \(\mathfrak{C}\) 中对类型化 \(\mathbb{T}\) 是可迁移的，其中 \(x_i (1 \leqslant i \leqslant n)\) 视为主基集， \(A_h (1 \leqslant h \leqslant m)\) 视为辅助基集：设 \(y_1, \ldots, y_n, f_1, \ldots, f_n\) 是两两不同的字母，并且不同于 \(x_i (1 \leqslant i \leqslant n)\) 、 \(s_j (1 \leqslant j \leqslant p)\) 、 \(\mathfrak{C}\) 的常量，以及 \(\mathbb{R}\) 或各项 \(A_h (1 \leqslant h \leqslant m)\) 中出现的所有字母；再设 \(\operatorname{Id}_h (1 \leqslant h \leqslant m)\) 表示 \(A_h\) 到自身的恒等映射。则关系

\begin{enumerate}
\def\labelenumi{(\arabic{enumi})}
\tightlist
\item
  “T \(\{x_{1},\ldots,x_{n},s_{1},\ldots,s_{p}\}\) 且 ( \(f_{1}\) 是 \(x_{1}\) 到 \(y_{1}\) 上的双射 ) 且 \ldots{} 且 ( \(f_{n}\) 是 \(x_{n}\) 到 \(y_{n}\) 上的双射 )''
\end{enumerate}

蕴含，在 \(\mathfrak{G}\) ，关系

\[
\mathrm{R} \left\{x _ {1}, \dots , x _ {n}, s _ {1}, \dots , s _ {p} \right\} \Longleftrightarrow \mathrm{R} \left\{y _ {1}, \dots , y _ {n}, s _ {1} ^ {\prime}, \dots , s _ {p} ^ {\prime} \right\},\tag{2}
\]

其中

\[
s _ {j} ^ {\prime} = \left\langle f _ {1}, \dots , f _ {n}, \operatorname{Id} _ {1}, \dots , \operatorname{Id} _ {m} \right\rangle^ {S_j} (s _ {j}) \quad (1 \leqslant j \leqslant p).\tag{3}
\]

在没有辅助集合的情况下，存在一个类似但更简单的定义。

例如，如果 \(n = p = 2\) 并且如果类型化 \(\mathbf{T}\) 是“\(s_1 \in x_1\) 并且 \(s_2 \in x_1\)''，关系 \(s_1 = s_2\) 是可迁移的。另一方面，关系 \(x_1 = x_2\) 不是可迁移的。

